% titel: Blatt 1, A5/A6: Schema der Materie / Teilchen-Tabelle
% Teil II: der durchgerechnete Weg.

\subsection*{A5 — Schema der Materie}
Das Schema ordnet die Materie (mit Vertretern):
\begin{center}
\begin{tabular}{ll}
Materie & \\
\quad Reine Stoffe & Elemente (O$_2$, Fe), Verbindungen (H$_2$O, CO$_2$) \\
\quad \\Gemische & homogen (Lösung, Legierung), heterogen (Suspension, Emulsion)
\end{tabular}
\end{center}
Die Kästchen im Blatt entsprechen: oben \emph{reine Stoffe} vs.\
\emph{Gemische}; darunter je Elemente/Verbindungen bzw.\ homogen/heterogen.

\subsection*{A6 — Teilchen-Tabelle}
Konkrete Einträge folgen der Definition aus A4: Spalte \emph{Protonen $Z$},
\emph{Neutronen $A-Z$}, \emph{Elektronen (neutral)}. Beispiel-Einträge:
O (8), P (15), Cu (29), Ag (47), U (92) — wie in A4 berechnet.

\begin{ergebnis}
Schema und Tabelle füllen sich über die Begriffe aus A4/A5; je reine
Stoffe/Gemische bzw.\ Z/N/E-Werte eintragen.
\end{ergebnis}

\begin{kontrolle}
Die Zuordnung ``reiner Stoff vs.\ Gemisch'' entscheidet die chemische
Zusammensetzung; die Z/N/E-Werte stammen aus den Nuklidzahlen.
\end{kontrolle}

\begin{merksatz}
Materie-Schema: reine Stoffe (Element/Verbindung) vs.\ Gemische
(homogen/heterogen). Teilchenzahlen aus der Nuklidbezeichnung.
\end{merksatz}
